\subsection{The complete reduced central kernel at the fixed finite base}

Retain the actual quotient $U$ of B.2.13, its normal
subgroup $N=L_0^6$, and the quotient
$E=U/N$ of order $2^{14}3^6$ in (B.S4).
Write $K_E$ for the inverse image of $N$ in $B$.
Thus $K\subseteq K_E$, and the marked quotient maps are
the maps induced by the same named generators $s,r$:
\[
\begin{gathered}
 U^{\mathrm c}=B/\overline{[B,K]}
       \longrightarrow E^{\mathrm c}
          =B/\overline{[B,K_E]},\\
 M=K/\overline{[B,K]}
       \longrightarrow M_E=K_E/\overline{[B,K_E]}.
\end{gathered}                                                   \tag{B.S32}
\]
Let $P_E=J/N\subseteq E$ be its actual pure subgroup.
For $t=s,r$, define
\[
 \kappa_{t,E}:C_{P_E}(t)\longrightarrow M_E,
       \qquad c\longmapsto[\widetilde c,t],                    \tag{B.S33}
\]
where the lift is taken in $E^{\mathrm c}$.  Centrality
of $M_E$ makes this independent of that lift.  The
commutator convention is the same convention as in
(B.S11); replacing all such commutators by their
inverses leaves every asserted image subgroup unchanged.

\begin{result}{B.2.20}{Proved}{The finite base cover and its complete torsion correction}
For these specified maps one has
\[
\begin{aligned}
 M_E&\cong C_2\times C_3,\\
 |\operatorname{im}\kappa_{s,E}|&=2,
       &|\operatorname{im}\kappa_{r,E}|&=3,\\
 M_E&=\operatorname{im}\kappa_{s,E}
                         \operatorname{im}\kappa_{r,E}.
\end{aligned}                                                   \tag{B.S34}
\]
The map $M\to M_E$ in (B.S32) is onto.  The maps
$C_J(s)\to C_{P_E}(s)$ and
$C_J(r)\to C_{P_E}(r)$ are also onto and intertwine
the respective $\kappa$ maps.  Consequently, for
$S=\operatorname{im}\kappa_s\operatorname{im}\kappa_r$
and $M_0=\ker(M\to M_E)$, the actual remaining reduced
kernel is
\[
\begin{gathered}
 M_0=\frac{K\cap\overline{[B,K_E]}}
                      {\overline{[B,K]}},\qquad
 M/S\cong M_0/(M_0\cap S).                                    \tag{B.S35}
\end{gathered}
\]
In particular the finite base calculation detects no
reduced obstruction; (B.S35) retains the entire kernel
which remains beyond that base.
\end{result}

\begin{proof}
The finite group $E$ is the actual marked group
constructed in B.2.13; its power--conjugate
presentation has twenty generators.  The following
complete cohomology and source-restriction calculation
uses that presentation.  With trivial coefficients
$\mathbb F_p$, the restrictions to the named torsion
factors and the subsequent central-cover test are
\[
\begin{array}{c|c|c|c|c|c}
 p&\dim H^2(E,\mathbb F_p)&
 \substack{\text{source}\\\text{restriction vector}}&
 \substack{\dim\text{ first}\\\text{marked layer}}&
 \dim H^2(E_p,\mathbb F_p)&
 \substack{\text{next}\\\text{constraint rank}}\\ \hline
 2&2&(0,1)&1&2&2\\
 3&2&(1,0)&1&1&1
\end{array}                                                     \tag{B.S36}
\]
Here $E_p\to E$ is the marked first central layer.
For $p=2$ the source restriction is the square of a
lift of $s$; for $p=3$ it is the cube of a lift of
$r$.  The other torsion generator has a unique
order-preserving central correction, since its order
is prime to $p$.  In the last column the constraints
are exactly the source-order relation together with
all commutators between a lift of the existing
central kernel and each power--conjugate generator.
Thus the last two columns exclude every further
central marked $C_p$ layer.  They impose centrality
over $E$, not merely centrality over $E_p$.

For completeness, this test determines the full
kernel rather than only its elementary quotients.
The canonical finite kernel is a finite abelian group
whose primes divide $|E|$.  Its maximal elementary
$p$ quotient has dimension one by the first three
columns.  Any larger $p$-primary kernel would have
a further quotient with kernel $C_p$ over this
first layer, with the whole kernel central over $E$;
that would give a nonzero class satisfying precisely
the last-column constraints.  Their full rank excludes
it.  Taking the marked fibre product of the two first
layers gives a group of order $6|E|$, generated by
the resulting $s,r$, with kernel $C_2\times C_3$.
It is therefore the full group $E^{\mathrm c}$.

In this group the exact centralizer calculation gives
\[
\begin{array}{c|c|c}
 t&|C_{P_E}(t)|&|\operatorname{im}\kappa_{t,E}|\\ \hline
 s&576&2\\
 r&72&3\\
 sr&288&1
\end{array}                                                     \tag{B.S37}
\]
The first two image orders prove (B.S34).  The last
row concerns the finite base only; the actual
centralizer of $\tau$ in $U$ is the one determined
in B.2.18.

All entries in (B.S36)--(B.S37), including the
higher central constraints, are computed in
\path{work/centralizer_code/check_universal_central_full.g}
and its included source files, with the complete
executed record in the matching \texttt{.out} file.
The calculation constructs each extension from its
cohomology cocycle and evaluates the displayed
relations as group words.  It also tests all twenty-four
actual target frames of B.2.12 in the full cover:
each has the unique central cycle correction of
B.2.6 and preserves both positive torsion conjugacy
classes.  A conjugator can be made pure by its named
torsion power, since the centralizer of the corresponding
permutation in $S_3$ is the corresponding cyclic
subgroup.  These are finite-base assertions throughout.

It remains to compare this certificate with $U$.
The equality $[B,L_0]=L_0$ of B.2.11, applied in each
of the six coordinates, gives $[J,N]=N$, hence
$[U,N]=N$.  Therefore the cokernel of $M\to M_E$ is
\[
 K_E/(K\overline{[B,K_E]})
                     \cong N/\overline{[U,N]}=1,
\]
proving surjectivity and the kernel formula in (B.S35).

Now let $t=s$ or $r$ and take $c\in C_{P_E}(t)$.
Choose a lift $h\in J$.  The two elements $hth^{-1}$
and $t$ have the same image in $E$ and the same
finite order.  Conjugation by $t$ permutes the six
factors of $N=L_0^6$ without fixed factors: it has
three two-cycles for $s$ and two three-cycles for $r$.
The factor maps around each cycle compose to the
identity, since $t$ has the stated order in $U$.
On such a permutation cycle every nonabelian
one-cocycle is a coboundary: choose its correcting
coordinate arbitrarily at the first factor and
solve successively around the cycle; the final
equation is exactly the finite-order cocycle relation.
Applying this independently on all cycles supplies
$n\in N$ such that $nh$ commutes with $t$.
It still maps to $c$, proving the two centralizer
surjections.  The commutator definitions and the
actual map (B.S32) give the asserted intertwining.

It follows from (B.S34) that $S$ maps onto $M_E$.
Every class in $M/S$ can consequently be represented
by an element of $M_0$, and its kernel there is
exactly $M_0\cap S$.  This proves the last isomorphism
in (B.S35), without identifying that remaining
quotient with the already computed finite base.
\end{proof}
